\documentclass[10pt,a4paper]{amsart}
\usepackage[margin=19mm]{geometry}
\usepackage{amsmath,amssymb,mathtools}
\usepackage[scr=rsfs,cal=euler]{mathalfa}
\usepackage{xfrac}
\usepackage[unicode,hidelinks,bookmarks=false]{hyperref}
\newcommand{\ie}{i.e.\ }
\newcommand{\cf}{cf.\ }
\newcommand{\resp}{respectively\ }
\newcommand{\Subsection}{\subsection}
\providecommand{\nobreakdash}{\nobreak\hbox{-}}
\setlength{\emergencystretch}{2em}
\multlinegap0pt
\usepackage{bm}
\usepackage{bbm}
\usepackage{dsfont}
\usepackage{xspace}
\usepackage{cancel}
\usepackage{mathtools}
\usepackage{amsthm}
\makeatletter
\@ifundefined{theorem}{\newtheorem{theorem}{Theorem}}{}
\@ifundefined{lemma}{\newtheorem{lemma}[theorem]{Lemma}}{}
\makeatother
\newcommand{\Z}{\mathbb Z}
\title[On Zero-sum Ramsey numbers of complete bipartite graphs]{On Zero-sum Ramsey numbers of complete bipartite graphs}
\author{Cheng Chi, Jialin He}
\date{Source: arXiv:2606.29216v2 (CC BY 4.0)}
\begin{document}
\maketitle

\noindent\emph{Source and licence.} On Zero-sum Ramsey numbers of complete bipartite graphs. Version arXiv:2606.29216v2 (CC BY 4.0), distributed under the Creative Commons Attribution 4.0 International licence, as recorded in the arXiv licence metadata for this exact version and shown on the abstract page https://arxiv.org/abs/2606.29216v2. Full rights notice: licenses/arxiv-2606.29216-CC-BY-4.0.txt.\par\smallskip

\noindent\textit{Editorial scope.} Counted unit: the Lemma whose source label is \texttt{lem:k23-structure} in the article (source0.tex lines 529--531, source label \texttt{lem:k23-structure}), delivered together with the article's own complete original proof of it (source0.tex lines 533--730, one \texttt{proof} environment of 198 lines). No mathematics is written, repaired or re-derived: statement and proof are the article's own text line for line. Required context retained uncounted: eq:k23-delta-type (lines 506--508); eq:k23-struct-degree (lines 516--518); eq:k23-struct-pair (lines 520--522); eq:k23-struct-triple (lines 524--526). Every definition, display and label of those lines is retained unchanged. Omitted from this package and not redistributed: 1--505, 509--515, 519--519, 523--523, 527--528, 532--532, 731--1028. Nothing inside a retained excerpt is omitted or altered: every retained body is delivered exactly as the article's lines are written, so no omission receipt of any kind accompanies this package. Citations: the retained excerpts carry no citation command at all, so no bibliography entry of the article is transcribed and no reference is redistributed. Reference adaptation: none was needed and none was made; every reference inside the delivered unit resolves inside it, so no unresolved reference or citation is produced. Printed slips kept literal: no printed word, symbol, hypothesis or number of the counted statement or of its proof was altered. Layout and class: the article's own class is \texttt{article} with packages including fontenc, inputenc, geometry, amsmath, amssymb, amsthm, mathtools, enumitem, xcolor, hyperref, cleveref, footmisc; the wrapper is the lane's pinned \texttt{amsart} editorial wrapper at 10pt on A4, which restates the theorem-like environments the retained text needs and adds the article's own macro definitions that the retained text needs (\texttt{\textbackslash Z}), with extra wrapper packages (bm, bbm, dsfont, xspace, cancel, mathtools). The delivered numbering is the unit's own. Assets: no retained span contains a figure environment, a graphics-inclusion command, a PDF-page inclusion, a table or a TikZ picture; no image file is copied into this package, the package declares no asset dependency and no image file is redistributed with it. Licence: version arXiv:2606.29216v2 of this article is distributed under the Creative Commons Attribution 4.0 International licence, as recorded in the arXiv licence metadata for this exact version and shown on the abstract page https://arxiv.org/abs/2606.29216v2. A full rights notice is delivered at licenses/arxiv-2606.29216-CC-BY-4.0.txt. Characters: every retained excerpt is pure ASCII, so the delivered engine, XeLaTeX with its default Latin Modern text font, reports no missing character.
\begin{equation}\label{eq:k23-delta-type}
    \delta_{yz}=0\text{ if }r_y\ne r_z\text{ and }\delta_{yz}\in\{1,2\}\text{ if }r_y=r_z.
\end{equation}

\begin{equation}\label{eq:k23-struct-degree}
    \sum_{z\in R\setminus\{y\}}\delta_{yz}=r_y+B.
\end{equation}

\begin{equation}\label{eq:k23-struct-pair}
    \delta_{yz}-b_y-b_z+r_y+r_z\ne0.
\end{equation}

\begin{equation}\label{eq:k23-struct-triple}
    \delta_{yz}+b_y+b_z-b_t-\delta_{yt}-\delta_{zt}\ne0.
\end{equation}

\begin{lemma}\label{lem:k23-structure}
    There is no finite set $X$ with data $r_x,b_x\in\Z_3$ and symmetric values $\delta_{xy}\in\Z_3$ for distinct $x,y\in X$ such that $|X|\ge 5$, $|X|\equiv2\pmod3$, $B\coloneqq\sum_xb_x=\sum_xr_x$, and conditions \eqref{eq:k23-delta-type}, \eqref{eq:k23-struct-degree}, \eqref{eq:k23-struct-pair}, and \eqref{eq:k23-struct-triple} all hold with $R$ replaced by $X$.
\end{lemma}

\begin{proof}
    For $i\in\Z_3$, let $X_i=\{x\in X:r_x=i\}$ and $C_i=\{b_x:x\in X_i\}$.

    First suppose that all three sets \(X_0,X_1,X_2\) are nonempty.
    If $i\ne j$, $x\in X_i$, and $y\in X_j$, then $\delta_{xy}=0$ by \eqref{eq:k23-delta-type}. Hence \eqref{eq:k23-struct-pair} gives
    \begin{equation}\label{eq:CiCj-avoid}
        b_x+b_y\ne i+j.
    \end{equation}
    Moreover, if \(x\in X_i\), \(y\in X_j\), and \(z\in X_\ell\), where \(i,j,\ell\) are distinct, then all three relevant \(\delta\)-terms are zero, and \eqref{eq:k23-struct-triple} gives
    \begin{equation}\label{eq:CiCjCl-avoid}
        b_z\ne b_x+b_y.
    \end{equation}
    These two restrictions imply $(C_0,C_1,C_2)=(\{0\},\{0\},\{1\})$ or $(C_0,C_1,C_2)=(\{0\},\{2\},\{0\})$.
    Indeed, if \(1\in C_0\), then \eqref{eq:CiCj-avoid} implies \(C_1\subseteq\{1,2\}\) and \(C_2\subseteq\{0,2\}\). Since \eqref{eq:CiCjCl-avoid} also forbids \(C_1+C_2\) from meeting \(C_0\), we obtain \(C_1=\{2\}\) and \(C_2=\{0\}\), which then violates \eqref{eq:CiCjCl-avoid}. Similarly, \(2\notin C_0\). Thus \(C_0=\{0\}\). Then \eqref{eq:CiCj-avoid} gives \(C_1\subseteq\{0,2\}\) and \(C_2\subseteq\{0,1\}\), while \eqref{eq:CiCjCl-avoid} implies \(C_1\cap C_2=\emptyset\) and forbids \(C_1+C_2\) from containing \(0\). The two possibilities displayed above are the only ones.

    Consider first
    \((C_0,C_1,C_2)=(\{0\},\{0\},\{1\}).\)
    For any internal edge of a layer \(X_i\), by~\eqref{eq:k23-delta-type}, the value of \(\delta\) is contained in \(\{1,2\}\). Applying \eqref{eq:k23-struct-triple} with two vertices in \(X_i\) and the third vertex in a suitable different layer implies that every internal \(\delta\)-edge in each \(X_i\) has value \(2\). Therefore \eqref{eq:k23-struct-degree} gives $2(|X_i|-1)=i+B$ for $i=0,1,2$. Here $B=|X_2|$ since the only nonzero \(b\)-values occur in \(X_2\).  These congruences imply $|X_0|\equiv0$, $|X_1|\equiv2$, and $|X_2|\equiv1$, so $|X|\equiv0$, contradicting $|X|\equiv2$.
    The second possibility,
    \((C_0,C_1,C_2)=(\{0\},\{2\},\{0\}),\)
    is analogous. In this case the same conditions imply that every internal \(\delta\)-edge has value \(1\). Thus $|X_i|-1=i+B$ for $i=0,1,2$. Since now \(B=2|X_1|\), these congruences imply $|X_0|\equiv0$, $|X_1|\equiv1$, and $|X_2|\equiv2$. Again \(|X|\equiv 0\), a contradiction. Hence at most two of the layers \(X_0,X_1,X_2\) are nonempty.

    Next suppose that exactly two layers are nonempty. If \(X_i\) and \(X_j\) are the two nonempty layers and \(u,v\in C_i\), then an internal \(\delta\)-edge in \(X_i\) joining vertices with \(b\)-values \(u\) and \(v\) must lie in
    \begin{equation}\label{eq:k23-D}
        \delta_i(u,v;C_j)
        =\{1,2\}
        \setminus (\{u+v-2i\}\cup \{c-u-v:c\in C_j\}).
    \end{equation}
    The first forbidden value comes from \eqref{eq:k23-struct-pair}, and the second set of forbidden values comes from \eqref{eq:k23-struct-triple} applied with the third vertex in the other layer $X_j$.

    We first claim that the two nonempty layers cannot be $X_1$ and $X_2.$
    Indeed, we first show that $|C_1|=|C_2|=1.$
    For cross-layer edges, \eqref{eq:CiCj-avoid} gives $u+v\ne0$ for all $u\in C_1$ and $v\in C_2$.
    Hence neither $C_1$ nor $C_2$ has three elements. Moreover, two two-element subsets of $\Z_3$ always contain a pair whose sum is $0$. Thus at least one of $C_1,C_2$ is a singleton. It remains to rule out the case in which the other has two elements.
    Suppose on the contrary that for some \(i\in\{1,2\}\), $|C_i|=2$ and $|C_j|=1,$ where $j=3-i.$
    Let $C_i=\{a,b\}$ and $C_j=\{c\}.$
    If $|X_j|\ge 2,$ then by~\eqref{eq:k23-D},
    the $\delta$-value of the edge of $X_j$ cannot be $2c-2j, a-2c, b-2c$. Since $a-2c\neq b-2c,$ one of $a-2c$ and $b-2c$ equals 0, i.e., $a=2c$ or $b=2c.$
    Then, by~\eqref{eq:CiCj-avoid}, $i+j\neq0,$ which is a contradiction to $j=3-i$.
    Therefore, $|X_j|=1,$ and thus $|X_i|\equiv1\pmod3$.
    Now $B=i|X_i|+j=i+j=0.$
    Choose the vertex of $X_j$, by~\eqref{eq:k23-struct-degree}, we get $0=j+B=j,$ which is a contradiction to $j=3-i$ and $i\in[2]$. Thus, $|C_1|=|C_2|=1.$

    Suppose that $0\in C_1,$ then by~\eqref{eq:CiCj-avoid}, $0\notin C_2$.
    If $C_1=\{0\}$ and $C_2=\{1\},$ then $B=|X_1|+2|X_2|=|X_2|,$ which implies that $|X|=|X_1|+|X_2|\equiv0\pmod3$, which contradicts the assumption.
    If $C_1=\{0\}$ and $C_2=\{2\},$ then $B=|X_1|+2|X_2|=2|X_2|,$ which implies that $|X_1|\equiv0, |X_2|\equiv2\pmod3$, since $|X|\equiv2\pmod3$.
    Choose two vertices from \(X_1\) and one vertex from $X_2$, by~\eqref{eq:k23-D}, we get a contradiction to the $\delta$ value of the edge in $X_1$.
    Therefore, we may assume that $0\notin C_1.$

    In the two cases \((C_1,C_2)=(\{1\},\{0\})\) and
    \((C_1,C_2)=(\{1\},\{1\})\), we have
    \(B\equiv2\), \(|X_1|\equiv2\), and \(|X_2|\equiv0\pmod3\).
    If \(C_2=\{0\}\), by~\eqref{eq:k23-D}, the $\delta$-value of edges of $X_1$ is 2, and if \(C_2=\{1\}\), then the $\delta$-value of edges of $X_1$ is 1.
    Choose a vertex from $X_1,$ by~\eqref{eq:k23-struct-degree}, if \(C_2=\{0\}\), then we get $2(|X_1|-1)=1+B=0,$
    if \(C_2=\{1\}\), then we get $|X_1|-1=1+B=0,$
    both of which contradict that $|X_1|\equiv2\pmod3$.
    Therefore, we may assume that $1\notin C_1.$

    For the case when $C_1=\{2\}$ and \(C_2=\{0\}\),
    we get that $B\equiv |X_1|+2|X_2|\equiv2|X_1|\pmod3$.
    Thus, $|X_1|\equiv2|X_2|\pmod3$, and hence $|X|\equiv0\pmod3$, contradicting $|X|\equiv2\pmod3$.
    If $C_1=\{2\}$ and \(C_2=\{2\}\), we get that $B\equiv |X_1|+2|X_2|\equiv2|X_1|+2|X_2|\pmod3$.
    Thus, $|X_1|\equiv0\pmod3$, $|X_2|\equiv2\pmod3$, and $B\equiv1\pmod3$.
    By~\eqref{eq:k23-D}, there is no allowable $\delta$-value for an internal edge in $X_1$. Since $X_1\ne\emptyset$, this gives $|X_1|=1$, contradicting $|X_1|\equiv0\pmod3$.
    Thus, the two nonempty layers can only be $X_0, X_1$ or $X_0, X_2.$

    Replacing \((r_x,b_x,\delta_{xy})\) by
    \((2r_x,2b_x,2\delta_{xy})\), if necessary, preserves all the conditions
    and interchanges the roles of \(X_1\) and \(X_2\). Hence it remains to
    consider the case \(X_0,X_1\ne\emptyset\).

    We first record a simple consequence of the pair and triple conditions.
    Suppose \(i\in\{0,1\}\), \(j=1-i\), and an internal edge in \(X_i\) joins
    vertices with \(b\)-values \(u,v\). If \(c\in C_j\), then, by
    \eqref{eq:k23-struct-pair} and by applying \eqref{eq:k23-struct-triple}
    with each of the three possible distinguished pairs, its \(\delta\)-value
    must lie in
    \[
        \Gamma_i(u,v;c)
        \coloneqq
        \{1,2\}\setminus
        \{u+v-2i,\ c-u-v,\ u+c-v,\ v+c-u\}.
    \]
    Thus \(\Gamma_i(u,v;c)\ne\emptyset\) whenever such an edge exists.

    For cross-layer edges, \eqref{eq:CiCj-avoid} gives $u+v\ne1$ for all $u\in C_0$ and $v\in C_1$.
    Hence neither \(C_0\) nor \(C_1\) has three elements; and if one of them
    has two elements, then the other is a singleton.
    First suppose that \(C_0=\{u,v\}\) and \(C_1=\{c\}\).
    The conditions \(u+c\ne1\) and \(v+c\ne1\) imply that
    \(c=2,0,1\) when \(C_0=\{0,1\},\{0,2\},\{1,2\}\), respectively.
    Substituting these three possibilities into \(\Gamma_0(u,v;c)\), only the first gives a nonempty allowable set.
    Thus the only possible case with \(|C_0|=2\) is
    \[
        (C_0,C_1)=(\{0,1\},\{2\}).
    \]
    If instead \(C_1=\{u,v\}\) and \(C_0=\{c\}\), the same check with \(\Gamma_1(u,v;c)\) gives an empty allowable set in all three cases.
    Therefore the only possible case in which one of \(C_0,C_1\) has two
    elements is $(C_0,C_1)=(\{0,1\},\{2\})$.

    It remains to consider the case where both \(C_0\) and \(C_1\) are
    singletons.  Write \(C_0=\{a\}\), \(C_1=\{b\}\), and let
    \(n_i=|X_i|\).  Then $an_0+bn_1\equiv B\equiv n_1$ and $n_0+n_1\equiv2\pmod3$.
    Together with \(a+b\ne1\), this leaves only
    \[
        (a,b)=(0,0),\ (0,2),\ (1,1),\ (2,1),
    \]
    since \((a,b)=(1,2)\) and \((2,0)\) would imply
    \(n_0+n_1\equiv0\pmod3\).

    The case \((a,b)=(1,1)\) is impossible: the two congruences give
    \(n_0\equiv0\) and \(n_1\equiv2\), so \(X_1\) contains an internal edge.
    However, we have $\Gamma_1(1,1;1)=\emptyset$, contradicting the existence of its \(\delta\)-value.

    Consequently the only remaining possibilities are
    \[
        (C_0,C_1)=(\{0\},\{0\}),\
        (\{0\},\{2\}),\
        (\{0,1\},\{2\}),\
        (\{2\},\{1\}).
    \]
    We now eliminate these four possibilities.
    First suppose
    \((C_0,C_1)=(\{0\},\{0\}).\)
    Then \(B=0\), and since \(B=\sum r_x\), we have \(|X_1|\equiv 0\). Hence \(|X_0|\equiv 2\). Inside \(X_0\), condition \eqref{eq:k23-struct-triple} implies that every triangle is monochromatic in the two possible \(\delta\)-values \(1\) and \(2\). Thus all internal \(\delta\)-edges of \(X_0\) have one common value \(c\in\{1,2\}\). The degree equation \eqref{eq:k23-struct-degree} at a vertex of \(X_0\) then gives
    \(c(|X_0|-1)=0,\)
    so \(|X_0|\equiv 1\), contradicting \(|X_0|\equiv 2\).
    Next suppose
    \((C_0,C_1)=(\{0\},\{2\}).\)
    Let \(n_i=|X_i|\). Then
    \(B=2n_1=n_1,\)
    so \(n_1\equiv 0\). The conditions imply that every internal \(\delta\)-edge of \(X_1\) has value \(1\). Therefore, by~\eqref{eq:k23-struct-degree}, a vertex of \(X_1\) gives
    \(n_1-1=1+B.\)
    Since \(B=0\), this gives \(n_1\equiv 2\), contradicting \(n_1\equiv 0\).
    Now suppose
    \((C_0,C_1)=(\{0,1\},\{2\}).\)
    Let
    \(n_{00}=|\{x\in X_0:b_x=0\}|\), \(n_{01}=|\{x\in X_0:b_x=1\}|\), \(n_1=|X_1|.\)
    The identity \(B=\sum r_x\) gives
    \(n_{01}+2n_1=n_1,\)
    so
    \(n_1\equiv 2n_{01}.\)
    Since \(|X|\equiv 2\), it follows that
    \(n_{00}\equiv 2.\)
    The conditions imply the following allowable internal \(\delta\)-values in \(X_0\):
    \[\delta=1 \text{ on } (0,0)\text{-pairs},\
        \delta=2 \text{ on } (0,1)\text{-pairs},\
        \delta=1 \text{ on } (1,1)\text{-pairs}.
    \]
    Thus the \(\delta\)-degree, inside \(X_0\), of a vertex with \(b\)-value \(0\) is
    \((n_{00}-1)+2n_{01}.\)
    By \eqref{eq:k23-struct-degree}, this must equal \(B=n_1\). Hence
    \((n_{00}-1)+2n_{01}=n_1.\)
    Using \(n_{00}\equiv 2\) and \(n_1\equiv 2n_{01}\), this becomes
    \(1+2n_{01}=2n_{01},\)
    a contradiction.
    Finally suppose
    \((C_0,C_1)=(\{2\},\{1\}).\)
    Let \(n_i=|X_i|\). Then
    \(B=2n_0+n_1=n_1,\)
    so \(n_0\equiv 0\). Since \(|X|\equiv 2\), we have \(n_1\equiv 2\). The conditions imply that every internal \(\delta\)-edge of \(X_0\) has value \(2\). Thus, by~\eqref{eq:k23-struct-degree}, a vertex of \(X_0\) gives
    \(2(n_0-1)=B=n_1,\) which yields \(1=2\), a contradiction. Hence exactly two layers cannot occur.

    It remains to consider the case in which all vertices lie in one layer.  Write $r_x=r$ for all $x\in X$.  Since $|X|\equiv2$, we have $B=\sum_xr_x=2r$.  Define $\beta_x=b_x+2r$.  Then
    \begin{equation}\label{eq:beta-sum-zero}
        \sum_x\beta_x=0.
    \end{equation}
    Condition \eqref{eq:k23-struct-pair} becomes $\delta_{xy}\ne\beta_x+\beta_y$.  Define $\eta_{xy}=\delta_{xy}+\beta_x+\beta_y$.  Then $\eta_{xy}=0$ unless $\beta_x+\beta_y=0$, and in the latter case $\eta_{xy}\in\{1,2\}$.  Equations \eqref{eq:k23-struct-degree} and \eqref{eq:k23-struct-triple} become
    \begin{equation}\label{eq:eta-degree}
        \sum_{y\ne x}\eta_{xy}=0\text{ for }x\in X,
    \end{equation}
    \begin{equation}\label{eq:eta-triple}
        \eta_{xy}-\eta_{xz}-\eta_{yz}+\beta_x+\beta_y+\beta_z+r\ne0
    \end{equation}
    for all distinct $x,y,z$.

    For \(i=0,1,2\), let $m_i=|\{x:\beta_x=i\}|$.  From $|X|\equiv2$ and \eqref{eq:beta-sum-zero}, we have
    \begin{equation}\label{eq:mi-basic}
        m_0+m_1+m_2\equiv2\text{ and }m_1+2m_2\equiv0.
    \end{equation}
    We first note that if \(m_0>0\), then
    \(m_0\equiv 1.\)
    Indeed, consider the vertices with \(\beta=0\). On this set, every \(\eta\)-edge has value \(1\) or \(2\), every vertex has \(\eta\)-degree \(0\) by \eqref{eq:eta-degree}, and \eqref{eq:eta-triple} becomes
    \(\eta_{xy}-\eta_{xz}-\eta_{yz}+r\ne 0.\)
    If \(r=0\), this implies that every triangle is monochromatic, so all edges have a common value \(c\in\{1,2\}\), and then \(c(m_0-1)=0\), giving \(m_0\equiv 1\). If \(r=1\), the allowed triangle patterns on the \(\beta=0\) vertices are exactly \((2,2,2)\) and \((1,1,2)\). Equivalently, the edges of value \(1\) form a cut. If this cut is nontrivial, say its parts have sizes $a$ and $m_0-a$ with $1\le a\le m_0-1$, then the \(\eta\)-degree of a vertex in the two parts gives $(m_0-a)+2(a-1)=0$ and $a+2(m_0-a-1)=0$, which are impossible in $\Z_3$. Hence all edges have value \(2\), and \(2(m_0-1)=0\), so \(m_0\equiv 1\). The case \(r=2\) is symmetric.

    Now consider the three possible values of \(r\).
    If \(r=0\), then three vertices with \(\beta=1\) would violate \eqref{eq:eta-triple}; hence \(m_1\le 2\). Similarly, \(m_2\le 2\). Together with \eqref{eq:mi-basic}, this leaves only
    \((m_1,m_2)=(0,0),\ (1,1),\ \text{or }(2,2).\)
    The case \((m_1,m_2)=(1,1)\) is impossible by \eqref{eq:eta-degree}, since the unique edge between the \(\beta=1\) vertex and the \(\beta=2\) vertex has nonzero \(\eta\)-value. If \((m_1,m_2)=(0,0)\), then \(m_0\equiv 2\), contradicting $m_0\equiv 1$. Thus \((m_1,m_2)=(2,2)\).
    Take two vertices \(u,v\) with \(\beta=1\) and one vertex \(z\) with \(\beta=2\). Applying \eqref{eq:eta-triple} with each of the three possible choices of the distinguished pair implies
    \(\eta_{uz}=\eta_{vz}=1.\)
    On the other hand, taking one vertex with \(\beta=1\) and two vertices with \(\beta=2\) implies that the corresponding \((1,2)\)-edges have \(\eta\)-value \(2\). This is impossible. Hence \(r\ne 0\).

    If \(r=1\), then one vertex with \(\beta=1\) and two vertices with \(\beta=2\) cannot occur. Indeed, if the two \((1,2)\)-edges have \(\eta\)-values $a,b\in\{1,2\}$ and the \((2,2)\)-edge has \(\eta=0\), then the three inequalities obtained from \eqref{eq:eta-triple} give $-a-b\ne0$, $a-b\ne0$ and $b-a\ne0$, which are impossible in $\Z_3$. Hence, if \(m_1>0\), then \(m_2\le 1\). If \(m_2=1\), then each vertex with \(\beta=1\) has exactly one possible nonzero incident \(\eta\)-edge, namely the edge to the unique \(\beta=2\) vertex, contradicting the zero-degree condition \eqref{eq:eta-degree}. Thus \(m_2=0\) whenever \(m_1>0\). Then, regardless of whether $m_1$ is equal to 0, \eqref{eq:mi-basic} gives \(m_1\equiv 0\), and hence \(m_0\equiv 2\), contradicting \(m_0\equiv 1\).
    The case \(r=2\) is symmetric to \(r=1\), obtained by multiplying all elements of \(\Z_3\) by \(2\).
    This final contradiction proves the lemma.
\end{proof}

\end{document}
